% !TeX program = xelatex
\documentclass[11pt,a4paper]{article}
\usepackage[margin=23mm]{geometry}
\usepackage{fontspec}
\setmainfont{Latin Modern Roman}
\usepackage{amsmath,amssymb,mathtools}
\usepackage{xcolor,enumitem,booktabs,longtable,array,xurl,hyperref}
\definecolor{muted}{HTML}{53636C}
\hypersetup{colorlinks=true,urlcolor=blue,linkcolor=blue}
\setlength{\parindent}{0pt}
\setlength{\parskip}{5pt}
\newcommand{\R}{\mathbb R}
\newcommand{\E}{\mathbb E}
\newcommand{\N}{\mathbb N}
\newcommand{\ind}{\mathbf 1}
\newcommand{\Var}{\operatorname{Var}}
\newcommand{\Cov}{\operatorname{Cov}}
\newcommand{\supp}{\operatorname{supp}}
\DeclareMathOperator*{\argmin}{arg\,min}
\DeclareMathOperator*{\argmax}{arg\,max}
\newcommand{\entrymeta}[3]{{\sffamily\small\color{muted}\textbf{\texttt{#1}}\quad|\quad #2\par #3\par}}
\newcommand{\Problemref}[1]{\texttt{#1}}
\newcommand{\JELref}[2]{\texttt{#2} (JEL \texttt{#1})}
\begin{document}
\section*{Strongly polynomial computation of an optimal finite contract}
\textbf{Problem ID:} \texttt{D86-21}\quad \textbf{JEL:} \texttt{D86}\par
% BEGIN ECONOMICS STATEMENT
\label{problem:OP-0596}


\label{audited:ECON-011}

{\sffamily\small\color{muted}\textbf{ECON-011} | P | Explicit published open question, at source date\par}

\subsection*{Definitions and formal target}

Let $n,m$ be positive integers, both variable parts of the input. A risk-neutral principal offers a payment vector $t\in\R_{\geq0}^{m}$, with payment $t_j$ contingent only on the observed outcome $j\in\{1,\ldots,m\}$. The agent's hidden action $a\in\{1,\ldots,n\}$ has cost $c_a\geq0$ and outcome probabilities $q_{aj}$; at least one action has zero cost. Principal rewards are $r_j\geq0$. The arrays $q,c,r$ are supplied explicitly as rational numbers in binary encoding and are known to both parties. Among actions maximizing the agent's expected payoff, ties are resolved by maximizing the principal's expected payoff; any remaining tie is immaterial to the objective.

\[
 \max_{\substack{a\in\{1,\ldots,n\}\\t\in\R_{\geq0}^{m}}}\sum_{j=1}^{m}q_{aj}(r_j-t_j)
 \quad\text{subject to}\quad
 \sum_{j=1}^{m}q_{aj}t_j-c_a\geq\sum_{j=1}^{m}q_{bj}t_j-c_b
 \quad\forall b\in\{1,\ldots,n\}.
\]
Does there exist an algorithm that, for every such input, returns an exactly optimal rational payment vector and an associated principal-favorable best-response action, using polynomially many arithmetic operations and comparisons in $n,m$, independently of numerical bit lengths, with intermediate encodings polynomial in the total binary input length?

\subsection*{Clarifications and required assumptions}

For every action $a$, $q_{aj}\geq0$ and $\sum_jq_{aj}=1$. Both parties maximize expected monetary payoff. Limited liability is $t_j\geq0$. The zero-cost action provides an outside option of at least zero because payments are nonnegative. Weak incentive inequalities are valid under the stated principal-favorable tie rule. For a candidate action $a$, infeasibility of its weak incentive inequalities means that this candidate is skipped; all original actions remain in the incentive comparisons for every other candidate. Feasibility here means implementability up to tie-breaking, rather than necessarily selection under the fixed tie rule for an arbitrary feasible payment vector. At a global optimum, the selected action can be taken to satisfy the principal-favorable tie rule, as in [S1, Proposition 3.1]. At least one candidate is feasible via the zero payment vector. For each feasible candidate, minimizing its expected payment is a rational linear program with objective bounded below by zero and an attained optimum; consequently the displayed global maximum is attained and has a rational optimal payment vector. ``Strongly polynomial'' bounds arithmetic-operation and comparison counts by a polynomial in $n,m$ and requires polynomial-size intermediate rational encodings; it does not mean total bit runtime ignores the cost of reading the input.

\subsection*{Variants and changes of scope}

\textbf{Regularity.} Fix a precise restriction on $q,c,r$ before asking for a restricted algorithm. \textbf{Fixed dimension.} Fix $n$ or $m$ independently of input length; this is another complexity regime. \textbf{Participation/ties.} A positive reservation utility or adversarial tie-breaking changes the program and may require strict inducement. \textbf{Large outcome space.} An implicitly encoded outcome space is not the explicit $n\times m$ input here. General strongly polynomial linear programming would imply a solution, but the source does not establish an equivalence.

\subsection*{Audit conclusion and source relationship}

The source explicitly asks the strongly polynomial question after stating ordinary polynomial-time solvability. The revised citation pins the checked 2024 preprint; it avoids treating ambiguous journal-year metadata as the date of that version.

\subsection*{References and exact source roles}

{\footnotesize\raggedright
\par\noindent\textbf{[S1]} Paul D\"utting, Michal Feldman, and Inbal Talgam-Cohen (2024 preprint). \emph{Algorithmic Contract Theory: A Survey}. arXiv:2412.16384v1, 20 December 2024. Journal DOI: 10.1561/0400000113; Foundations and Trends in Theoretical Computer Science 16(3--4), 211--412. The date here identifies the checked preprint, without conflating journal-year conventions.\par \textit{Verified locator / source role:} Section 3.1, computational discussion immediately after Observation 3.3.\par \url{https://arxiv.org/html/2412.16384v1}\par\smallskip
}

\subsection*{Common conventions: Additional-source status and mathematical conventions}

\label{audited:conventions}

Of the 100 supplied ECON entries, 65 distinct problems appear here; 32 close
matches refine earlier entries, and three duplicate questions map to existing
entries. Appendix~\ref{app:audited-crosswalk} lists every destination.
The attachment's P/R/S/U distinctions and source audits are retained. Its
precise mathematical formalizations are editorial unless the individual audit
says otherwise. Candidate subsidy targets ECON-004--006 retain their U flags;
the resolved approval-core question ECON-007 updates OP-0202 above.
The following supplied conventions govern both this part and the attributed
ECON refinements elsewhere in the catalogue, subject to local restrictions.

\textbf{Parameterized questions.} A problem family lists inputs that must be fixed before an instance is evaluated: population, horizon, policies, information, data law, model class and welfare weights. These are required choices, not quantities the citations are assumed to specify. Undefined applications should not be treated as identified estimands.

\textbf{Indices and integrability.} Sets of agents, firms and regions are finite unless a population measure is explicitly used. \(i,j\) index units, \(r\) regions, \(t\) dates and \(h\) horizons. \(\mathbb E\) and \(\Pr\) refer to the declared population and law; \(\mathbf1\{A\}\) is an indicator. Every displayed expectation and discounted sum needs existence in the stated domain. Infinite horizons use a discount in \((0,1)\) and a convergence condition; finite horizons may use discount one.

\textbf{Optimization and existence.} A displayed \(\max\), \(\min\), or \(\arg\max\) is conditional on attainment. For finite feasible menus this is immediate; general spaces need continuity and compactness/coercivity or another existence argument. Without attainment, the target is the supremum/infimum and arbitrarily near-optimal feasible choices. \(\inf\varnothing=+\infty\). Empty feasibility and an unidentified target are different issues.

\textbf{Potential outcomes.} \(Y_i(z)\) is the outcome under a specified intervention, not the observed conditional mean among people choosing \(z\). In binary comparisons \(\Delta Y_i=Y_i(1)-Y_i(0)\). Specify assignment, treatment versions, compliance, information and observation horizon. Interference is allowed under a full policy vector; compressed exposure notation needs a justified mapping. No interference, consistency and overlap are substantive assumptions, not automatic consequences of notation.

\textbf{Populations and observation.} Fix baseline eligibility and population weights. Conditioning on post-policy employment, survival, relocation or program uptake can change the population being compared. Death is not ordinary loss to follow-up; outcomes truncated by death need a composite convention, joint survival target, or explicitly defined survivor stratum. Fertility-changing interventions need a population functional rather than an unexplained fixed descendant cohort. Measurement and missingness models must be supplied.

\textbf{Identification.} A structural model \(m\in\mathcal M\) implies observable law \(P_m\) and target \(\theta(m)\). Identification means
\[
 P_{m_1}=P_{m_2}\ \Longrightarrow\ \theta(m_1)=\theta(m_2)
 \quad\text{for all }m_1,m_2\in\mathcal M .
\]
Without identification, the target is
\[
 \Theta(P;\mathcal M)=\{\theta(m):m\in\mathcal M,\ P_m=P\}.
\]
Writing \(\theta(P)\) before establishing identification can hide incompatible latent models. Confidence sets additionally need a sampling law, confidence level, asymptotic sequence and uniformity class. Point identification, coverage and informative precision are separate properties.

\textbf{Mediation.} Total effects of randomized assignment are distinct from cross-world natural indirect effects. Belief/effort-channel effects require a meaningful mediator intervention and additional measurement, exclusion or mediation assumptions. Randomization of the initial policy alone does not supply those assumptions. Report bounds or interventional alternatives when the stronger target is unsupported.

\textbf{Equilibrium.} A counterfactual \(W(z)\), \(p(z)\), or \(\mu_t^z\) requires individual optimization, feasibility, government/resource budgets, market clearing and state-transition laws. State equilibrium existence and selection, or report ranges over compatible equilibria. Initial-state integration includes subsequent endogenous transitions; current-state integration must use the policy-dependent distribution. Reduced-form invariance after policy is not assumed.

\textbf{Welfare accounts.} Scores, utility, health, dollars and ecological quantities require explicit conversion before addition. Fees, taxes, subsidies, wages and extortion payments are transfers between parties; distributional weights can make them welfare relevant, but their face value is not automatically a resource loss. State ownership, geographic boundaries, foreign-income weights and fiscal finance. Do not add profits to household welfare already including dividends, or count land capitalization and the underlying benefit twice. A benefit-cost expression must distinguish gross benefits from net welfare and subtract every real cost exactly once.

\textbf{Derivatives and risk.} Log derivatives require positive arguments; local responses require admissible perturbations and specified equilibrium branches. Differentiation under expectations needs regularity/domination, and boundaries may allow only one-sided derivatives. Risk uses a specified law; ambiguity uses a declared nonempty law set on a common history space. Adaptive policies must be nonanticipative. A robust criterion is an editorial choice unless attributed explicitly.

\textbf{Scope overlaps.} Allocation, collective choice, contracts, household bargaining and coordinated policy overlap economics with optimization and game theory. They are retained because they appeared in the original catalogue; no claim of a strictly game-theory-free selection is made.
% END ECONOMICS STATEMENT
\end{document}
